% Shared by both studies' REPORT.tex. The estimator does not differ between
% them, so its description should not either -- only the pipeline around it
% does, and each report writes its own "This analysis" subsection after this.
%
% Supplies \subsection* only; the including document owns \section*{Methods}.

\subsection*{The model}

bipred is bivariate LDpred: one joint fit of two traits that share a single LD
reference. Each variant is assigned to one of four latent states --- causal for
neither trait, for trait 1 only, for trait 2 only, or for \emph{both} --- with
probabilities $(\pi_{00}, \pi_{10}, \pi_{01}, \pi_{11})$. A trait-1-only effect
is drawn from $N(0, s_1)$, a trait-2-only effect from $N(0, s_2)$, and a
both-causal pair from $N(0, \Sigma)$ with
$\Sigma = \left[\begin{smallmatrix} s_1 & s_{12} \\ s_{12} & s_2
\end{smallmatrix}\right]$. The off-diagonal $s_{12}$ is the effect covariance
inside the shared component, and is the only place the two traits couple.

The indicator is \emph{per trait}, not a single shared one. Whether the two
traits' causal variants co-occur is therefore learned as $\pi_{11}$ rather than
assumed: correlated traits borrow strength through the both-causal component,
and disjoint traits can drive $\pi_{11}$ toward zero so the two fits decouple.
That is what makes the polygenic-overlap summary an estimate rather than a
restatement of the prior.

Fitting is Gibbs sampling over the LD blocks. Each sweep evaluates the four
bivariate-Gaussian likelihoods of the residual marginal estimate at a variant,
samples that variant's state, then draws its effects; the mixture weights $\pi$
and the variance components $(s_1, s_2, s_{12})$ are re-estimated every sweep
under a damping factor of $0.2$. Reported quantities are posterior means over
the retained iterations: per-trait SNP heritability \texttt{h2}, the genetic
correlation \texttt{rg}, posterior-mean effects, and a MiXeR-style overlap
summary --- \texttt{n\_causal} per trait, \texttt{n\_shared},
\texttt{frac\_shared}, and the correlation \texttt{rho\_beta} of effects within
the shared component.

\subsection*{What rho\_beta is, and is not}

\texttt{rg} is a genome-wide average: genetic covariance divided by the
geometric mean of the two heritabilities, over every variant in the fit.
\texttt{rho\_beta} is the effect correlation \emph{conditional on a variant
being causal for both traits}. The two answer different questions, and they come
apart exactly when one trait's heritability sits in very few loci. Neither is a
causal estimate.

\texttt{frac\_shared} is $\pi_{11}$ divided by the smaller of the two marginal
causal fractions, so it asks what proportion of the less polygenic trait's
causal set is also causal for the other.

\subsection*{Sample overlap}

Overlapping samples correlate the two traits' sampling errors. The fit accepts
that correlation as \texttt{cross\_corr}, set per pair from the cross-trait LD
Score regression intercept, used as-is and never clipped into range. An
intercept can also reflect correlated confounding rather than shared samples, so
this absorbs overlap without identifying it.

\subsection*{The LD-consistency screen}

Per-variant filters --- minor allele frequency, imputation quality, a
chi-square cap, sample size --- judge each variant in isolation, so none of them
can see a variant disagreeing with the variants it is correlated with. The
screen tests exactly that. Within a window the variants are split at random in
half and each $z$-score in one half is predicted from the other through the LD,
\[
  \hat{z}_a = R_{aB}\,R_{BB}^{+}\,z_B,
  \qquad
  T_a = \frac{(z_a - \hat{z}_a)^2}{1 - R_{aB}\,R_{BB}^{+}\,R_{Ba}}
  \;\sim\; \chi^2_1
\]
under consistency. Variants far from what their neighbours predict are dropped,
and the split is repeated so each variant is tested from several directions.

It runs against the same numeric LD representation the fit will use, so the
screen and the sampler see the same matrix. It is DENTIST-inspired, not DENTIST:
window construction, partition schedule, eigenvalue regularisation and removal
policy all differ from the published pipeline (Chen et al.\ 2021,
\emph{Nature Communications} 12:7117), so that pipeline's calibration does not
transfer.

\subsection*{Divergence diagnostics}

A bivariate fit can diverge silently. Given summary statistics that disagree
with the LD reference, the sampler places large opposing effects on variants in
near-perfect LD, and these cancel to a plausible-looking \texttt{h2}. Three
diagnostics are recorded, and any one of them raises the warning:
$\sum \beta^2 / h^2$ above $10$ (the cancellation ratio, reported in both
studies), $\max|\beta|$ above $25$ slab standard deviations, and retained-trace
drift above $1.25$. All three were calibrated on one real LDL $\times$ CAD fit
that diverged and the same fit after screening repaired it, so each separates
the two regimes by more than an order of magnitude rather than by a hair.

\textbf{An estimate carrying the divergence warning is not an estimate of
anything}, and is reported only to show the contrast.

\subsection*{LD reference and harmonisation}

All fits use the European UK Biobank HapMap3 LD reference, built with ldpred3's
bigsnpr converter at the pinned revision \texttt{5d86ac9d}. Summary statistics
are harmonised to it by rsID with allele matching, flipping the effect sign
where the effect and reference alleles are swapped; strand-ambiguous and
unmatched variants are dropped. Odds-ratio releases are taken to logs before
harmonisation, which preserves the sign across a flip. Case/control traits are
fitted at an effective $N$ of $4/(1/n_{\text{case}} + 1/n_{\text{ctrl}})$. BLAS
is pinned to one thread throughout.
